\noindent\textbf{Resultado.} \[x_1=e^t(C_1\cos2t+C_2\sin2t),\qquad x_2=e^t(C_2\cos2t-C_1\sin2t),\quad C_1,C_2\in\mathbb R.\]

\noindent\textit{Desarrollo.} De la primera fila, $x_2=(x_1'-x_1)/2$. Al diferenciar y sustituir en la segunda,
\[\frac{x_1''-x_1'}2=-2x_1+\frac{x_1'-x_1}2,\qquad x_1''-2x_1'+5x_1=0.\]
Las raíces características son $1\pm2i$, de modo que $x_1$ tiene la forma indicada. La reconstrucción $x_2=(x_1'-x_1)/2$ da la segunda componente. Recíprocamente, esa definición satisface la primera fila y convierte la segunda exactamente en la ecuación escalar; no se añaden ni pierden soluciones. Los datos $(x_1,x_2)$ y $(x_1,x_1')$ se relacionan por una matriz invertible, por lo que las dos constantes son independientes. Las soluciones son globales y su factor $e^t$ explica el crecimiento de la amplitud, con frecuencia angular $2$.
